\en{We will start with Kummer's solution (Thm.~\ref{\en{EN}omegastrich}) and use Clausen's Formula (Thm.~\ref{\en{EN}satzclausen}) and the Fourier series (Thm.~\ref{\en{EN}fouriertheorem}) to prove the Main Theorem~\ref{\en{EN}hauptformel}.
This chapter follows the paper \cite{\en{EN}glebov} of Chen and Glebov.}%
\de{Wir werden, ausgehend von Kummers Lösung (Thm.~\ref{\en{EN}omegastrich}) und mit Hilfe der Clausen-Formel (Thm.~\ref{\en{EN}satzclausen}) und den Fourierdarstellungen (Thm.~\ref{\en{EN}fouriertheorem}), das Haupttheorem~\ref{\en{EN}hauptformel} beweisen.
Dieses Kapitel orientiert sich stark am Paper von Chen und Glebov \cite{\en{EN}glebov}.}

\en{First we give an overview over the names of the basic periods and basic quasiperiods of the equivalent lattices $L_\tau$, $L_J$ and $\tilde L$ (cf.~Def.~\ref{\en{EN}definLJ} and Def.~\ref{\en{EN}defltilde}). And throughout this chapter, we assume that $L=\mathbb Z \omega_1 + \mathbb Z \omega_2$ equals $L_\tau$ with $\omega_1=1$:}%
\de{Zunächst stellen wir die Bezeichnungen für die Perioden und Quasiperioden der drei äquivalenten Gitter $L_\tau$, $L_J$ und $\tilde L$ zusammen (vgl.~Def.~\ref{\en{EN}definLJ} und Def.~\ref{\en{EN}defltilde}), wobei wir in diesem Kapitel voraussetzen, dass $L=\mathbb Z \omega_1 + \mathbb Z \omega_2$ von der Form $L_\tau$ mit $\omega_1=1$ ist:}

{\vspace*{1mm}\hspace*{2mm}\renewcommand{\arraystretch}{2}%
    \begin{tabular}{l|l}\label{\en{EN}tabgitter}
    \textbf{ \en{Lattices and Basic Periods}\de{Gitter und Perioden}} & \textbf{ \en{Basic Quasiperiods}\de{Quasiperioden} }\\
    \hline
    ~$\displaystyle L_\tau: (\omega_1,\omega_2)=(1,\tau)$
    & ~$\displaystyle (\eta_1,\eta_2)=(\eta_1(L_\tau),\eta_2(L_\tau))$\\
    ~$\displaystyle L_J: (\Omega_1,\Omega_2)=\sqrt{\frac{g_3(\tau)}{g_2(\tau)}}\cdot(1,\tau)$~
    &~$\displaystyle (H_1,H_2)=\sqrt{\frac{g_2(\tau)}{g_3(\tau)}}\cdot(\eta_1(L_\tau),\eta_2(L_\tau))$\\
    ~$\displaystyle \tilde L: (\tilde\omega_1,\tilde\omega_2)=\Delta(\tau)^{\frac 1 {12}}\cdot(1,\tau)$~
    & ~$\displaystyle (\tilde\eta_1,\tilde\eta_2)=\Delta(\tau)^{-\frac 1 {12}}\cdot(\eta_1(L_\tau),\eta_2(L_\tau))$\\
    \end{tabular}\vspace*{2mm}}

\en{If we use one of the expressions $\eta_1$, $\eta_2$, $g_2$, $g_3$ or $\Delta$ in this chapter, it shall denote $\eta_1(L_\tau)$, $\eta_2(L_\tau)$, $g_2(\tau)$, $g_3(\tau)$ or $\Delta(\tau)$.
And (like in Remark~\ref{\en{EN}bemwurzel}) we won't care about the choice of the branch of the roots in the intermediate steps, only when we arrive at the Main Theorem~\ref{\en{EN}hauptformel}.}%
\de{Wenn wir im Folgenden $\eta_1$, $\eta_2$, $g_2$, $g_3$ oder $\Delta$ schreiben, meinen wir damit $\eta_1(L_\tau)$, $\eta_2(L_\tau)$, $g_2(\tau)$, $g_3(\tau)$ bzw.~$\Delta(\tau)$.
Außerdem werden wir uns wie in Bem.~\ref{\en{EN}bemwurzel} bei den Zwischenschritten keine Gedanken zur Wahl der Wurzel machen, erst im Haupttheorem~\ref{\en{EN}hauptformel}.}

\begin{thm}\label{\en{EN}thm35}%(Gleichung (4.5) in \cite{\en{EN}chuddekker}):
\en{For $k=1$ and for $k=2$ it holds:}%
\de{Für $k=1$ und für $k=2$ gilt:}
$$\tilde\eta_k = -\sqrt{12}J^\frac 2 3\sqrt{J-1}\cdot\frac{d\tilde\omega_k}{dJ}$$
\end{thm}
\begin{proof}
\en{For $\sqrt{\frac{g_3}{g_2}}$, we calculated the representation~(\ref{\en{EN}mue}) on page~\pageref{\en{EN}mue}:}%
\de{Für $\sqrt{\frac{g_3}{g_2}}$ haben wir auf S.~\pageref{\en{EN}mue} die Darstellung~(\ref{\en{EN}mue}) berechnet:}
\begin{align}
    \sqrt{\frac{g_3}{g_2}}&=\underbrace{J^{-\frac 1 6} \cdot \left(\frac{J-1}{27}\right)^{\frac 1 4}}_{=:A(J)} \cdot~\Delta^\frac 1 {12}\label{\en{EN}g2g3}
\end{align}
\en{This yields a direct connection between $L_J$ and $\tilde L$:}%
\de{Hieraus folgt ein direkter Zusammenhang zwischen $L_J$ und $\tilde L$:}
\begin{align*}
    (\Omega_1,\Omega_2)&=\sqrt{\frac{g_3}{g_2}}\cdot(1,\tau)=\sqrt{\frac{g_3}{g_2}}\cdot\Delta^{-\frac 1{12}}\cdot(\tilde\omega_1,\tilde\omega_2) = A(J) \cdot(\tilde\omega_1,\tilde\omega_2)\\
    (H_1,H_2)&=\sqrt{\frac{g_2}{g_3}}\cdot(\eta_1,\eta_2)=\sqrt{\frac{g_2}{g_3}}\cdot\Delta^{\frac 1{12}}\cdot(\tilde\eta_1,\tilde\eta_2) = \frac{1}{A(J)} \cdot(\tilde\eta_1,\tilde\eta_2)
\end{align*}
\en{We derive the first line by $J$ and obtain:}%
\de{Die erste Zeile leiten wir noch nach $J$ ab und erhalten:}
\begin{align*}
    \frac{d\Omega}{dJ} &= A(J) \cdot\frac{d\tilde\omega}{dJ} + \left(-\frac{1}{6J}+\frac{1}{4(J-1)}\right)\cdot A(J) \cdot\tilde\omega\\
    &=A(J) \cdot \left(\frac{d\tilde\omega}{dJ} + \frac{J+2}{12J(J-1)}\cdot\tilde\omega\right)
\end{align*}
\en{Now we use this new equation in~(\ref{\en{EN}dOdJ}) on p.~\pageref{\en{EN}dOdJ} and get the desired equation:}%
\de{Wir setzen die soeben gefundenen Zusammenhänge nun auf Seite~\pageref{\en{EN}dOdJ} in Gleichung~(\ref{\en{EN}dOdJ}) ein und erhalten die zu beweisende Gleichung:}\belowdisplayskip=-12pt
\begin{align*}
    36 J(J-1)\overbrace{A(J)\cdot \left(\frac{d\tilde\omega}{dJ} + \frac{J+2}{12J(J-1)}\cdot\tilde\omega\right)}^{\frac{d\Omega}{dJ}} &= 3(J+2)\overbrace{A(J)\tilde\omega}^\Omega -2(J-1)\overbrace{\frac{1}{A(J)}\tilde\eta}^H\\
    \Longrightarrow\qquad \frac{2(J-1)}{A(J)}\cdot\tilde\eta &= -36J(J-1)A(J)\cdot\frac{d\tilde\omega}{dJ}\\
    \Longrightarrow\qquad \tilde\eta
    &=-\sqrt{12}J^\frac 2 3\sqrt{J-1}\cdot\frac{d\tilde\omega}{dJ}
\end{align*}\end{proof}

\begin{defi}\label{\en{EN}defis2}
\en{We denote the following non-holomorphic function by $s_2$:}%
\de{Die folgende nicht-holomorphe Funktion nennen wir $s_2$:}
$$s_2(\tau):=\frac{E_4(\tau)}{E_6(\tau)}\cdot E_2^*(\tau)\qquad\text{\en{with}\de{mit}}\qquad E_2^*(\tau) := E_2(\tau)-\frac{3}{\pi \im(\tau)}$$
\en{where $E_k(\tau)$ are the normalized Eisenstein series from Thm.~\ref{\en{EN}fouriertheorem}.}%
\de{wobei $E_k(\tau)$ die normierten Eisensteinreihen aus Thm.~\ref{\en{EN}fouriertheorem} sind.}
\end{defi}

\begin{bem}
\en{$s_2$ is an example of an "almost holomorphic modular form". These are functions in $\mathbb H$ which transform like a modular form -- i.e.~the values of $s_2$ are equal for equivalent lattices -- but instead of being holomorphic, they are polynomials in $\frac{1}{\im(\tau)}$ with holomorphic coefficients.
And in Prop.~\ref{\en{EN}exts2rat} we will see that certain values of $s_2(\tau)$ are rational.}%
\de{$s_2$ ist eine \glqq fast holomorphe Modulform\grqq. Diese sind Funktionen in $\mathbb H$, die wie eine Modulform transformieren -- auch für $s_2$ gilt, dass es bei äquivalenten Gittern den gleichen Wert annimmt -- aber sie sind nicht holomorph, sondern Polynome in $\frac{1}{\im(\tau)}$ mit holomorphen Koeffizienten.
Und in Satz~\ref{\en{EN}exts2rat} werden wir sehen, dass gewisse Werte von $s_2(\tau)$ rational sind.}
\end{bem}

\begin{thm}\label{\en{EN}thmglg10}
\en{It holds}\de{Es gilt} $$ \eta_1-\frac{3g_3}{2g_2}s_2(\tau) = \frac \pi {\im(\tau)} $$
\end{thm}
\begin{proof}
\en{We change the representations of $\eta_1$, $g_2$ and $g_3$ from Thm.~\ref{\en{EN}fouriertheorem} into representations of $E_2$, $E_4$ and $E_6$.
Then we put these into the definition of $s_2$ and obtain:}%
\de{Wir lösen die Darstellungen von $\eta_1$, $g_2$ und $g_3$ aus Theorem~\ref{\en{EN}fouriertheorem} nach $E_2$, $E_4$ und $E_6$ auf, setzen die Ergebnisse in die Definition von $s_2(\tau)$ ein und erhalten:}\belowdisplayskip=-12pt
\begin{align*}
    s_2(\tau) = \frac{\frac{3g_2}{4\pi^4}}{\frac{27g_3}{8\pi^6}}\left(\frac{3\eta_1}{\pi^2}-\frac{3}{\pi \im(\tau)}\right)
    &=\frac{2\pi^2}{9}\cdot\frac{g_2}{g_3}\left(\frac{3\eta_1}{\pi^2}-\frac{3}{\pi \im(\tau)}\right)\\
    &= \frac{2g_2}{3g_3}\cdot\eta_1-\frac{2\pi g_2}{3g_3 \im(\tau)}\\
    \Longrightarrow\qquad \frac{2g_2}{
    3g_3}\cdot\eta_1 - s_2(\tau) &= \frac{2\pi g_2}{3g_3 \im(\tau)}\quad\left|\cdot\frac{3g_3}{2 g_2}\right.\\
    \Longrightarrow\qquad \eta_1 - \frac{3g_3}{2 g_2}\cdot s_2(\tau) &= \frac{\pi}{\im(\tau)}
\end{align*}
\end{proof}


\begin{thm}\label{\en{EN}thm42}
\en{For all $\tau$ with $\im(\tau)>1\ko25$ it holds}%
\de{Für alle $\tau$ mit $\im(\tau)>1\ko25$ gilt}
$$\frac{1}{2\pi \im(\tau)}\sqrt{\frac{J}{J-1}} = \frac{1-s_2(\tau)} 6 F^2 - J \frac d {dJ} F^2$$
\en{where we denote}\de{Dabei ist} $F = {_2F_1}\left(\frac{1}{12},\frac{5}{12};1;\frac 1 {J}\right)$ \en{and}\de{und} $J=J(\tau)$.
\end{thm}

\begin{proof}
\en{First we will prove Eq.~(\ref{\en{EN}etaeins}) and~(\ref{\en{EN}gdreigzwei}). Then we will combine them with Prop.~\ref{\en{EN}thmglg10}.}%
\de{Wir werden Glg.~(\ref{\en{EN}etaeins}) und~(\ref{\en{EN}gdreigzwei}) beweisen und sie dann in Satz~\ref{\en{EN}thmglg10} einsetzen.}

\en{Thm.~\ref{\en{EN}omegastrich} tells us that $\tilde\omega_1 = \Delta^{\frac 1 {12}} = \frac{2\pi}{\sqrt[4]{12}}\cdot J^{-\frac{1}{12}}\cdot F$ holds in the given region. We derive this by $J$ and obtain (using the product rule):}%
\de{Zunächst sagt Theorem~\ref{\en{EN}omegastrich}, dass $\tilde\omega_1 = \Delta^{\frac 1 {12}} = \frac{2\pi}{\sqrt[4]{12}}\cdot J^{-\frac{1}{12}}\cdot F$ im angegebenen Bereich gilt. Die Ableitung dieser Gleichung nach $J$ liefert mit der Produktregel:}
$$\frac{d\tilde\omega_1}{dJ}=\frac{2\pi}{\sqrt[4]{12}}\cdot J^{-\frac 1 {12}} \cdot \left(\frac{-1}{12J}\cdot F + \frac{dF}{dJ}  \right)$$
\en{Using this in the equation of Prop.~\ref{\en{EN}thm35} we get:}%
\de{Das setzen wir nun in die Gleichung aus Satz~\ref{\en{EN}thm35} ein:}
\begin{align*}
\tilde\eta_1 &= -\sqrt{12}J^\frac 2 3\sqrt{J-1}\cdot\overbrace{\frac{2\pi}{\sqrt[4]{12}}\cdot J^{-\frac 1 {12}} \cdot \left(\frac{-1}{12J}\cdot F + \frac{dF}{dJ}  \right)}^{\frac{d\tilde\omega_1}{dJ}}\nonumber\\
&= -2\pi\sqrt[4]{12}J^\frac 7{12} \sqrt{J-1} \cdot \left(\frac{-1}{12J}\cdot F + \frac{dF}{dJ}  \right)
\end{align*}
\en{From the table on p.~\pageref{\en{EN}tabgitter} we conclude $\eta_1 = \tilde\eta_1\cdot\Delta^{\frac{1}{12}}$.
Here we can use the newly found representation of $\tilde \eta_1$ and the representation of $\Delta^{\frac{1}{12}}$ from Thm.~\ref{\en{EN}omegastrich}. This yields:}%
\de{Aus der Tabelle auf S.~\pageref{\en{EN}tabgitter} entnehmen wir $\eta_1 = \tilde\eta_1\cdot\Delta^{\frac{1}{12}}$. Hier setzen wir die eben gefundene Darstellung von $\tilde \eta_1$ sowie die Darstellung von $\Delta^{\frac{1}{12}}$ aus Thm.~\ref{\en{EN}omegastrich} ein und erhalten:}
\begin{align}
    \eta_1 = \tilde\eta_1\cdot\Delta^{\frac{1}{12}} &= - 2\pi\sqrt[4]{12}J^\frac 7{12} \sqrt{J-1} \cdot \left(\frac{-1}{12J}\cdot F + \frac{dF}{dJ}  \right)\cdot \frac{2\pi}{\sqrt[4]{12}}\cdot J^{-\frac{1}{12}}\cdot F \nonumber\\
    &=\frac{\pi^2}{3}\cdot \sqrt{\frac{J-1}{J}}\cdot F^2  -  2\pi^2\sqrt{J(J-1)}\cdot \underbrace{2F\frac{dF}{dJ}}_{\frac{d}{dJ}\lk F^2\rk}\label{\en{EN}etaeins}
\end{align}
\en{Next we take eq.~(\ref{\en{EN}g2g3}) from p.~\pageref{\en{EN}g2g3} and use the representation of $\Delta^{\frac{1}{12}}$ from Thm.~\ref{\en{EN}omegastrich}:}%
\de{Nun nehmen wir Glg.~(\ref{\en{EN}g2g3}) von S.~\pageref{\en{EN}g2g3} und setzen die Darstellung von $\Delta^{\frac{1}{12}}$ aus Thm.~\ref{\en{EN}omegastrich} ein:}
\begin{align}
    \frac{3g_3}{2g_2} &= \frac 3 2 \cdot \left(J^{-\frac 1 6} \cdot \left(\frac{J-1}{27}\right)^{\frac 1 4} \cdot~\Delta^\frac 1 {12}\right)^2
    = \frac 3 2 \cdot J^{-\frac 1 3} \cdot \left(\frac{J-1}{27}\right)^{\frac 1 2} \cdot \Delta^\frac 1 {6}\nonumber\\
    &= \frac 3 2 \cdot J^{-\frac 1 3} \cdot \frac{\sqrt{J-1}}{\sqrt{27}}\cdot \frac{4\pi^2}{\sqrt{12}}\cdot J^{-\frac{1}{6}}\cdot F^2 = \frac{\pi^2}{3}\cdot\sqrt{\frac{J-1}{J}}\cdot F^2\label{\en{EN}gdreigzwei}
\end{align}
\en{Now we use~(\ref{\en{EN}etaeins}) and~(\ref{\en{EN}gdreigzwei}) in the equation of Prop.~\ref{\en{EN}thmglg10}:}%
\de{Nun setzen wir~(\ref{\en{EN}etaeins}) und~(\ref{\en{EN}gdreigzwei}) in die Gleichung aus Satz~\ref{\en{EN}thmglg10} ein:}
\begin{align*}
    \overbrace{\frac{\pi^2}{3}\cdot \sqrt{\frac{J-1}{J}}\cdot F^2  -  2\pi^2\sqrt{J(J-1)}\cdot \frac{d}{dJ}\lk F^2\rk}^{\eta_1} - \overbrace{\frac{\pi^2}{3}\cdot\sqrt{\frac{J-1}{J}}\cdot F^2 }^{\frac{3g_3}{2g_2}}\cdot~s_2(\tau) &= \frac \pi {\im(\tau)}\\
    \Longrightarrow\quad \frac{\pi^2}{3}\cdot \sqrt{\frac{J-1}{J}}\cdot F^2 \cdot \left(1-s_2(\tau)\right) - 2\pi^2\sqrt{J(J-1)}\cdot \frac{d}{dJ}\lk F^2\rk &= \frac \pi {\im(\tau)}
\end{align*}
\en{Finally we multiply with $\sqrt{\frac{J}{J-1}}\cdot\frac{1}{2\pi^2}$ and obtain:}%
\de{Schließlich multiplizieren wir das mit $\sqrt{\frac{J}{J-1}}\cdot\frac{1}{2\pi^2}$ und erhalten:}
\begin{align*}
    \frac{1-s_2(\tau)}{6}\cdot F^2 - J\cdot \frac{d}{dJ}\lk F^2\rk &= \sqrt{\frac{J}{J-1}}\cdot\frac{1}{2\pi\im(\tau)}
\end{align*}
\en{This finishes the proof of Prop.~\ref{\en{EN}thm42}.}%
\de{Somit ist Satz~\ref{\en{EN}thm42} bewiesen.}
\end{proof}


\begin{thm}\label{\en{EN}darst}
\en{For the square of the following hypergeometric function it holds:}%
\de{Für das Quadrat der folgenden hypergeometrischen Funktion gilt:}
\begin{align*}
    \left({_2F_1}\lk\frac{1}{12},\frac{5}{12};1;z\rk\right)^2 &= \sum_{n=0}^\infty \frac{(6n)!}{(3n)!(n!)^3} \frac{z^n}{12^{3n}}
\end{align*}
\end{thm}
\begin{proof}
  \en{From Clausen's formula (Thm.~\ref{\en{EN}satzclausen} on p.~\pageref{\en{EN}satzclausen}) we get using Def.~\ref{\en{EN}defihyp} on p.~\pageref{\en{EN}defihyp}:}%
  \de{Aus der Formel von Clausen (Theorem~\ref{\en{EN}satzclausen} auf Seite~\pageref{\en{EN}satzclausen}) folgt mit Hilfe von Definition~\ref{\en{EN}defihyp} auf Seite~\pageref{\en{EN}defihyp}:}
  \begin{align}
      \left({_2F_1}\lk\frac{1}{12},\frac{5}{12};1;z\rk\right)^2 = {_3F_2}\lk\frac{1}{6},\frac{5}{6},\frac{1}{2};1,1;z\rk = \sum_{n=0}^{\infty} \frac{\left(\frac{1}{6}\right)_n\cdot\left(\frac{5}{6}\right)_n\cdot\left(\frac{1}{2}\right)_n}{(1)_n \cdot (1)_n}\cdot\frac{z^n}{n!}\label{\en{EN}asdf}
  \end{align}
  \en{Since $(1)_n=n!$ we only have to express $\left(\frac{1}{6}\right)_n \cdot \left(\frac{5}{6}\right)_n \cdot \left(\frac{1}{2}\right)_n$ in terms of factorials:}%
  \de{Es gilt aber $(1)_n=n!$, also müssen wir nur noch $\left(\frac{1}{6}\right)_n \cdot \left(\frac{5}{6}\right)_n \cdot \left(\frac{1}{2}\right)_n$ vereinfachen:}
  
    \en{If $a=\frac p q$ is the quotient of two natural numbers, the Def.~\ref{\en{EN}defpoch} of the Pochhammer symbols yields:}%
    \de{Wenn $a=\frac p q$ ein Verhältnis zweier natürlicher Zahlen ist, dann gilt aufgrund der Definition~\ref{\en{EN}defpoch} der Pochhammer-Symbole:}
    $$\left(\frac p q\right)_n = \prod_{k=1}^n\left(\frac p q + k - 1\right) = q^{-n} \prod_{k=1}^n\left(p + kq - q\right).$$
    \en{This yields (cf.~\cite[Lemma 4.1]{\en{EN}glebov}):}%
    \de{Hieraus folgt (vgl.~\cite[Lemma 4.1]{\en{EN}glebov}):}
    \begin{align*}
         \left(\frac{1}{6}\right)_n \cdot \left(\frac{5}{6}\right)_n \cdot \left(\frac{3}{6}\right)_n
         &= 6^{-3n}\prod_{k=1}^n(6k-5)(6k-3)(6k-1)\\
         &= 6^{-3n}\cdot 1\cdot 3\cdot 5 \cdot 7 \cdots (6n-1)\\
         &= 6^{-3n}\cdot\frac{(6n)!}{2\cdot 4\cdot 6\cdots 6n}\\
         &= 6^{-3n}\cdot\frac{(6n)!}{2^{3n}\cdot (3n)!}= \frac{(6n)!}{(3n)!\cdot 12^{3n}}
    \end{align*}
    \en{If we put this along with $(1)_n=n!$ into eq.~(\ref{\en{EN}asdf}) we get:}%
    \de{Wenn wir das in Gleichung~(\ref{\en{EN}asdf}) einsetzen, erhalten wir mit $(1)_n=n!$:}
    \begin{align*}
      \left({_2F_1}\lk\frac{1}{12},\frac{5}{12};1;z\rk\right)^2  = \sum_{n=0}^{\infty} \frac{(6n)!}{(3n)!\cdot 12^{3n}\cdot n! \cdot n!}\cdot\frac{z^n}{n!}
  \end{align*}
  \en{which concludes the proof.}%
  \de{und das ist was zu zeigen war.}
\end{proof}

\pagebreak
\begin{theo}[\en{Main Theorem}\de{Haupttheorem}]\label{\en{EN}hauptformel}
\en{For all $\tau$ with $\im(\tau)>1\ko25$ we have the following identity due to David and Gregory Chudnovsky, first published in 1988 \cite[Eq.~(1.4)]{\en{EN}chud1988}:}%
\de{Für alle $\tau$ mit $\im(\tau)>1\ko25$ gilt die folgende Gleichung von David und Gregory Chudnovsky aus dem Jahr 1988 \cite[Glg.~(1.4)]{\en{EN}chud1988}:}
\begin{align*}
    \frac{1}{2\pi \im(\tau)}\sqrt{\frac{J(\tau)}{J(\tau)-1}} &= \sum_{n=0}^\infty \left( \frac{1-s_2(\tau)}{6} + n \right)\cdot \frac{(6n)!}{(3n)!(n!)^3}\cdot \frac{1}{\left(1728J(\tau)\right)^n}
\end{align*}
\en{Here $\sqrt{\phantom{J}}$ denotes the principal branch of the square root.}%
\de{Hierbei bezeichnet $\sqrt{\phantom{J}}$ den Hauptzweig der Quadratwurzel.}
\end{theo}

\begin{proof}
\en{For the proof we combine the differential equation from Prop.~\ref{\en{EN}thm42} with the representation from Prop.~\ref{\en{EN}darst}:}%
\de{Für den Beweis kombinieren wir die Differentialgleichung aus Satz~\ref{\en{EN}thm42} mit der Darstellung aus Satz~\ref{\en{EN}darst}:}

\en{First we use the same notation like in Prop.~\ref{\en{EN}thm42} and denote the function $F(J)={_2F_1}\lk\frac{1}{12},\frac{5}{12};1;\frac 1 {J}\rk$ with $J=J(\tau)$.
Then we call $G(z)=\left({_2F_1}\lk\frac{1}{12},\frac{5}{12};1;z\rk\right)^2$.
Then it holds $\left(F(J)\right)^2=G(z)$ with $z=\frac 1 J$.
In Prop.~\ref{\en{EN}thm42} we proved: $$\frac{1}{2\pi \im(\tau)}\sqrt{\frac{J}{J-1}} = \frac{1-s_2(\tau)} 6 F^2 - J \frac d {dJ} F^2$$}%
\de{Wir bezeichnen zunächst wie in Satz~\ref{\en{EN}thm42} die Funktion $F(J)={_2F_1}\lk\frac{1}{12},\frac{5}{12};1;\frac 1 {J}\rk$ mit $J=J(\tau)$.
Darüberhinaus nennen wir $G(z)=\left({_2F_1}\lk\frac{1}{12},\frac{5}{12};1;z\rk\right)^2$.
Dann haben wir mit $z=\frac 1 J$ nämlich $\left(F(J)\right)^2=G(z)$.
In Satz~\ref{\en{EN}thm42} haben wir bereits $$\frac{1}{2\pi \im(\tau)}\sqrt{\frac{J}{J-1}} = \frac{1-s_2(\tau)} 6 F^2 - J \frac d {dJ} F^2$$
bewiesen.}
\en{Now we transform this differential equation of $F(J)$ into one of $G(z)$, like in the proof of Prop.~\ref{\en{EN}satzbj}.
From $J=\frac{1}{z}$ we get $\frac{dJ}{dz}=\frac{-1}{z^2}$ and $\frac{dz}{dJ}=-z^2$. This yields:}%
\de{Wir wandeln diese Differentialgleichung für $F(J)$ in eine für $G(z)$ um, ähnlich wie im Beweis von Satz~\ref{\en{EN}satzbj}.
Aus $J=\frac{1}{z}$ folgt $\frac{dJ}{dz}=\frac{-1}{z^2}$ und $\frac{dz}{dJ}=-z^2$. Also gilt}
$$J \frac d {dJ} (F(J))^2 = \frac 1 z \cdot \frac {dz}{dJ} \cdot \frac{d}{dz} G(z) = -z \cdot \frac{d}{dz} G(z)$$
\en{This produces a differential equation of $G(z)$:}%
\de{Dies liefert uns eine Differentialgleichung für $G(z)$:}
$$\frac{1}{2\pi \im(\tau)}\sqrt{\frac{J}{J-1}} = \frac{1-s_2(\tau)} 6 G(z) + z \frac d {dz} G(z)\qquad\text{\en{with}\de{mit} }J=\frac 1 z$$
\en{Here we use the representation of $G(z)$ from Prop.~\ref{\en{EN}darst}:}%
\de{Hier können wir noch die Darstellung von $G(z)$ aus Satz~\ref{\en{EN}darst} einsetzen:}
\begin{align*}
    G(z)&=\sum_{n=0}^\infty \frac{(6n)!}{(3n)!(n!)^3} \frac{z^n}{12^{3n}}\\
    \Longrightarrow\quad z \frac d {dz} G(z) &= \sum_{n=0}^\infty \frac{(6n)!}{(3n)!(n!)^3} \frac{n\cdot z^n}{12^{3n}}\\
    \Longrightarrow\quad \frac{1}{2\pi \im(\tau)}\sqrt{\frac{J}{J-1}} &= \sum_{n=0}^\infty \left(\frac{1-s_2(\tau)} 6 + n\right)\cdot \frac{(6n)!}{(3n)!(n!)^3}\cdot \frac{z^n}{12^{3n}}
\end{align*}
\en{And finally we get the statement of Thm.~\ref{\en{EN}hauptformel}, if we use $z=\frac 1 J$ and $12^3=1728$.}%
\de{Und wir erhalten schließlich mit $z=\frac 1 J$ und $12^3=1728$ die zu beweisende Aussage.}

\en{In the intermediate steps until here we never cared about the choice of the root (cf.~Remark~\ref{\en{EN}bemwurzel}).
Our result in Thm.~\ref{\en{EN}hauptformel} is thus (until now) only proven correct up to a root of unity.
Thus we still have to prove for any $\tau$ with $\im(\tau)>1\ko25$ that our result doesn't contain any hidden complex roots of unity:}%
\de{In allen Zwischenschritten bis hier war uns die Wahl des Zweigs der Wurzeln egal (vgl.~Bemerkung~\ref{\en{EN}bemwurzel}),
deshalb haben wir bisher auch nur bewiesen, dass die Gleichung aus Thm.~\ref{\en{EN}hauptformel} korrekt bis auf eine komplexe Einheitswurzel ist.
Wir müssen also noch beweisen, dass die Gleichung für irgendein $\tau$ mit $\im(\tau)>1\ko25$ exakt ist und keine versteckten Einheitswurzeln enthält:}

\en{We choose $\tau_8=\frac{i\sqrt{8}}{2}=i\sqrt{2}$.
Here we get that $q=e^{2\pi i \tau_8} = e^{-2\pi\sqrt{2}}$ is a real number.
Thus (from Thm.~\ref{\en{EN}fouriertheorem}) both $J(\tau_8)$ and $s_2(\tau_8)$ are also real-valued.
The approximations together with the error estimates from Thm.~\ref{\en{EN}theonaeherJ} and~\ref{\en{EN}theonaehers2} tell us that both $J(\tau_8)$ and $\frac{1-s_2(\tau_8)}{6}$ are \emph{positive} real numbers.
This shows that all quantities in the equation of Thm.~\ref{\en{EN}hauptformel} are real-valued and positive at $\tau=\tau_8$.


This proves that if we choose the principal branch of the square root (which is positive on the positive real axis), the equation is exact without any hidden roots of unity at $\tau=\tau_8$.
The region $\im(\tau)>1\ko25$ is connected, both sides of the equation depend continuously on $\tau$ and are not zero.
This proves that the equation is exact without any hidden roots of unity for all $\tau$ with $\im(\tau)>1\ko25$, if the principal branch of the square root is chosen.}%
\de{Wir wählen $\tau_8=\frac{i\sqrt{8}}{2}=i\sqrt{2}$.
Hier gilt, dass $q=e^{2\pi i \tau_8} = e^{-2\pi\sqrt{2}}$ eine reelle Zahl ist.
Die Darstellungen in Thm.~\ref{\en{EN}fouriertheorem} zeigen dann, dass $J(\tau_8)$ und $s_2(\tau_8)$ ebenfalls reellwertig sind.
Weiter folgt aus den Näherungen und den Fehlerabschätzungen aus Thm.~\ref{\en{EN}theonaeherJ} und~\ref{\en{EN}theonaehers2}, dass sowohl $J(\tau_8)$ und $\frac{1-s_2(\tau_8)}{6}$ \emph{positive} reelle Zahlen sind.
Somit ist bewiesen, dass für $\tau=\tau_8$ alle Größen in der Gleichung des Haupttheorems~\ref{\en{EN}hauptformel} reellwertig und positiv sind.


Das zeigt, dass in der Gleichung bei $\tau=\tau_8$ keine komplexen Einheitswurzeln vergessen wurden, sofern man den Hauptzweig der Quadratwurzel wählt, welcher den postiven Radikanden positive Werte zuordnet.
Das Gebiet $\im(\tau)>1\ko25$ ist zusammenhängend, beide Seiten der Gleichung hängen stetig von $\tau$ ab und sind nicht Null.
Hieraus folgt, dass die Gleichung aus Thm.~\ref{\en{EN}hauptformel} auch für alle anderen $\tau$ mit $\im(\tau)>1\ko25$ exakt erfüllt ist, sofern man wie gefordert den Hauptzweig der Quadratwurzel wählt.}
\end{proof}
